\section{AMC 10 10A 2012}
        \begin{enumerate}
        	\item (\textit{AMC 10 10A 2012}) Cagney can frost a cupcake every 20 seconds and Lacey can frost a cupcake every $30$ seconds. Working together, how many cupcakes can they frost in $5$ minutes?


 $\textbf{(A)}\ 10\qquad\textbf{(B)}\ 15\qquad\textbf{(C)}\ 20\qquad\textbf{(D)}\ 25\qquad\textbf{(E)}\ 30$


 
 \textbf{Answer: }D



\textbf{Solution 1}

Cagney can frost one in $20$ seconds, and Lacey can frost one in $30$ seconds. Working together, they can frost one in $\frac{20\cdot30}{20+30} = \frac{600}{50} = 12$ seconds. In $300$ seconds ($5$ minutes), they can frost $\boxed{\textbf{(D)}\ 25}$ cupcakes.

 



\textbf{Solution 2}

In $300$ seconds ($5$ minutes), Cagney will frost $\dfrac{300}{20} = 15$ cupcakes, and Lacey will frost $\dfrac{300}{30} = 10$ cupcakes. Therefore, working together they will frost $15 + 10 = \boxed{\textbf{(D)}\ 25}$ cupcakes.

 



\textbf{Solution 3}

Since Cagney frosts $3$ cupcakes a minute, and Lacey frosts $2$ cupcakes a minute, they together frost $3+2=5$ cupcakes a minute. Therefore, in $5$ minutes, they frost $5\times5 = 25 \Rightarrow \boxed{\textbf{(D)}}$

 


	\item (\textit{AMC 10 10A 2012}) A square with side length $8$ is cut in half, creating two congruent rectangles. What are the dimensions of one of these rectangles?


 $\textbf{(A)}\ 2\ \text{by}\ 4\qquad\textbf{(B)}\ \ 2\ \text{by}\ 6\qquad\textbf{(C)}\ \ 2\ \text{by}\ 8\qquad\textbf{(D)}\ 4\ \text{by}\ 4\qquad\textbf{(E)}\ 4\ \text{by}\ 8$


 
 \textbf{Answer: }E



\textbf{Solution}

Cutting the square in half will bisect one pair of sides while the other side will remain unchanged. Thus, the new square is $\frac{8}{2}$ by $8$, or  $\boxed{\textbf{(E)}\ 4\ \text{by}\ 8}$.

 


	\item (\textit{AMC 10 10A 2012}) A bug crawls along a number line, starting at $-2$. It crawls to $-6$, then turns around and crawls to $5$. How many units does the bug crawl altogether?


 $\textbf{(A)}\ 9\qquad\textbf{(B)}\ 11\qquad\textbf{(C)}\ 13\qquad\textbf{(D)}\ 14\qquad\textbf{(E)}\ 15$


 
 \textbf{Answer: }E



\textbf{Solution}


                    \begin{figure}[H]
                        \centering
                        \includegraphics[width=0.5\textwidth]{Solution/AMC_10_10A_2012_3_Solution_Figure_1.png}
                    \end{figure}
                    

 Crawling from $-2$ to $-6$ takes it a distance of $4$ units. Crawling from $-6$ to $5$ takes it a distance of $11$ units. Add $4$ and $11$ to get $\boxed{\textbf{(E)}\ 15}$

 


	\item (\textit{AMC 10 10A 2012}) Let $\angle ABC = 24^\circ$ and $\angle ABD = 20^\circ$. What is the smallest possible degree measure for angle $CBD$?


 $\textbf{(A)}\ 0\qquad\textbf{(B)}\ 2\qquad\textbf{(C)}\ 4\qquad\textbf{(D)}\ 6\qquad\textbf{(E)}\ 12$


 
 \textbf{Answer: }C



\textbf{Solution}

$\angle ABD$ and $\angle ABC$ share ray $AB$. In order to minimize the value of $\angle CBD$, $D$ should be located between $A$ and $C$. 

 $\angle ABC = \angle ABD + \angle CBD$, so $\angle CBD = 4$. The answer is $\boxed{\textbf{(C)}\ 4}$

 


	\item (\textit{AMC 10 10A 2012}) Last year $100$ adult cats, half of whom were female, were brought into the Smallville Animal Shelter. Half of the adult female cats were accompanied by a litter of kittens. The average number of kittens per litter was $4$. What was the total number of cats and kittens received by the shelter last year?


 $\textbf{(A)}\ 150\qquad\textbf{(B)}\ 200\qquad\textbf{(C)}\ 250\qquad\textbf{(D)}\ 300\qquad\textbf{(E)}\ 400$


 
 \textbf{Answer: }B



\textbf{Solution}

Half of the 100 adult cats are female, so there are $\frac{100}{2}$ = $50$ female cats.  Half of those female adult cats have a litter of kittens, so there would be $\frac{50}{2}$ = $25$ litters.  Since the average number of kittens per litter is 4, this implies that there are $25 \times 4$ = $100$ kittens.  So the total number of cats and kittens would be $100 + 100$ = $\boxed{\textbf{(B)}\ 200}$

 


	\item (\textit{AMC 10 10A 2012}) The product of two positive numbers is $9$. The reciprocal of one of these numbers is $4$ times the reciprocal of the other number. What is the sum of the two numbers?


 $\textbf{(A)}\ \frac{10}{3}\qquad\textbf{(B)}\ \frac{20}{3}\qquad\textbf{(C)}\ 7\qquad\textbf{(D)}\ \frac{15}{2}\qquad\textbf{(E)}\ 8$


 
 \textbf{Answer: }D



\textbf{Solution}

Let the two numbers equal $x$ and $y$. From the information given in the problem, two equations can be written:

 $xy=9$

 $\frac{1}{x}=4 \left( \frac{1}{y} \right)$

 Therefore, $4x=y$

 Replacing $y$ with $4x$ in the equation,

 $4x^2=9$

 So $x=\frac{3}{2}$ and $y$ would then be $4 \times$ $\frac{3}{2}=6$

 The sum would be $\frac{3}{2}+6$ = $\boxed{\textbf{(D)}\ \frac{15}{2}}$

 



\textbf{Solution 2}

Set up your first equation which is $xy=9$ and your second being $\dfrac{1}{x} = \dfrac4y$. Then, in the first equation, rearrange it to become $x=\dfrac9y$. Now plug this in to your second equation, and you should get $\dfrac{y}{9}=\dfrac{4}{y}$. Cross multiply to get $y^2=36$, and simplify to get $y=6$. Notice how the problem mentioned that they were positive integers, so we don't consider $-6$. We plug $6$ back into our first equation to get $6x=9$ which makes $x$ come out to be $1.5$. We add $x+y=6+1.5=7.5$. Notice answer choice $D$ is the only fraction that simplifies to $1.5$, so answer option $\boxed{D}$ must be the answer.

 


	\item (\textit{AMC 10 10A 2012}) In a bag of marbles, $\frac{3}{5}$ of the marbles are blue and the rest are red. If the number of red marbles is doubled and the number of blue marbles stays the same, what fraction of the marbles will be red?


 $\textbf{(A)}\ \frac{2}{5}\qquad\textbf{(B)}\ \frac{3}{7}\qquad\textbf{(C)}\ \frac{4}{7}\qquad\textbf{(D)}\ \frac{3}{5}\qquad\textbf{(E)}\ \frac{4}{5}$


 
 \textbf{Answer: }C



\textbf{Solution 1}

Assume that there are 5 total marbles in the bag. The actual number does not matter, since all we care about is the ratios, and the only operation performed on the marbles in the bag is doubling.

 There are 3 blue marbles in the bag and 2 red marbles. If you double the amount of red marbles, there will still be 3 blue marbles but now there will be 4 red marbles. Thus, the answer is $\boxed{\textbf{(C)}\ \frac{4}{7}}$.

 



\textbf{Solution 2}

Let us say that there are $x$ marbles in the bag. Therefore, $\frac{3x}{5x}$ are blue, and $\frac{2x}{5x}$ are red. When the red marbles are doubled, we now have $\frac{2*2x}{5x+2x} = \frac{4x}{7x} = \frac{4}{7} \Rightarrow \boxed{\textbf{(C)}}$

 


 



\textbf{Solution 3}

Suppose there are $5$ marbles in total. Since $\frac{3}{5}$ are blue, there are $3$ blue marbles and $2$ red marbles.

 If the number of red marbles is doubled, then the number of red marbles becomes $4$, while the number of blue marbles remains $3$.

 Thus the new total number of marbles is $3+4=7$, and the fraction of marbles that are red is

 $\frac{4}{7}$.

 Therefore, the answer is $\boxed{\frac{4}{7}}$.

 $\text{(Solution by Trashboycan)}$

 


	\item (\textit{AMC 10 10A 2012}) The sums of three whole numbers taken in pairs are $12$, $17$, and $19$. What is the middle number?


 $\textbf{(A)}\ 4\qquad\textbf{(B)}\ 5\qquad\textbf{(C)}\ 6\qquad\textbf{(D)}\ 7\qquad\textbf{(E)}\ 8$


 
 \textbf{Answer: }D



\textbf{Solution 1}

Let the three numbers be equal to $a$, $b$, and $c$. We can now write three equations:

 $a+b=12$

 $b+c=17$

 $a+c=19$

 Adding these equations together, we get that

 $2(a+b+c)=48$ and

 $a+b+c=24$

 Substituting the original equations into this one, we find

 $c+12=24$

 $a+17=24$

 $b+19=24$

 Therefore, our numbers are 12, 7, and 5. The middle number is $\boxed{\textbf{(D)}\ 7}$

 



\textbf{Solution 2 (Faster)}

Let the three numbers be $a$, $b$ and $c$ and $a<b<c$. We get the three equations: 

 $a+b=12$

 $a+c=17$

 $b+c=19$

 To isolate $b$, We add the first and last equations and then subtract the second one.

 $(a+b)+(b+c)-(a+c) = 12+19-17 \Rightarrow 2b=14 \Rightarrow b = 7$

 Because $b$ is the middle number, the middle number is $\boxed{\textbf{(D)}\ 7}$

 


	\item (\textit{AMC 10 10A 2012}) A pair of six-sided dice are labeled so that one die has only even numbers (two each of $2$, $4$, and $6$), and the other die has only odd numbers (two each of $1$, $3$, and $5$). The pair of dice is rolled. What is the probability that the sum of the numbers on the tops of the two dice is $7$?


 $\textbf{(A)}\ \frac{1}{6}\qquad\textbf{(B)}\ \frac{1}{5}\qquad\textbf{(C)}\ \frac{1}{4}\qquad\textbf{(D)}\ \frac{1}{3}\qquad\textbf{(E)}\ \frac{1}{2}$


 
 \textbf{Answer: }D



\textbf{Solution}

The total number of combinations when rolling two dice is $6*6 = 36$.

 There are three ways that a sum of 7 can be rolled. $2+5$, $4+3$, and $6+1$. There are two 2's on one die and two 5's on the other, so there are a total of 4 ways to roll the combination of 2 and 5. There are two 4's on one die and two 3's on the other, so there are a total of 4 ways to roll the combination of 4 and 3. There are two 6's on one die and two 1's on the other, so there are a total of 4 ways to roll the combination of 6 and 1. Add $4 + 4 + 4 = 12$.

 Thus, our probability is $\frac{12}{36} = \frac{1}{3}$. The answer is $\boxed{\textbf{(D)}\ \frac13}$.

 



\textbf{Solution 2}

Assume we roll the die with only evens first. For whatever value rolled, there are exactly 2 faces on the odd die that makes the sum 7. The odd die has 6 faces, so our probability is $\boxed{\textbf{(D)}\ \frac13}$.

 


	\item (\textit{AMC 10 10A 2012}) Mary divides a circle into $12$ sectors. The central angles of these sectors, measured in degrees, are all integers and they form an arithmetic sequence. What is the degree measure of the smallest possible sector angle?


 $\textbf{(A)}\ 5\qquad\textbf{(B)}\ 6\qquad\textbf{(C)}\ 8\qquad\textbf{(D)}\ 10\qquad\textbf{(E)}\ 12$


 
 \textbf{Answer: }C



\textbf{Solution 1}

Let $a_1$ be the first term of the arithmetic progression and $a_{12}$ be the last term of the arithmetic progression. From the formula of the sum of an arithmetic progression (or arithmetic series), we have $12*\frac{a_1+a_{12}}{2}=360$, which leads us to $a_1 + a_{12} = 60$. $a_{12}$, the largest term of the progression, can also be expressed as $a_1+11d$, where $d$ is the common difference. Since each angle measure must be an integer, $d$ must also be an integer. We can isolate $d$ by subtracting $a_1$ from $a_{12}$ like so: $a_{12}-a_1=a_1+11d-a_1=11d$. Since $d$ is an integer, the difference between the first and last terms, $11d$, must be divisible by $11.$ Since the total difference must be less than $60$, we can start checking multiples of $11$ less than $60$ for the total difference between $a_1$ and $a_{12}$. We start with the largest multiple, because the maximum difference will result in the minimum value of the first term. If  the difference is $55$, $a_1=\frac{60-55}{2}=2.5$, which is not an integer, nor is it one of the five options given. If the difference is $44$, $a_1=\frac{60-44}{2}$, or $\boxed{\textbf{(C)}\ 8}$

 -Solution by Rhiju

 



\textbf{Solution 2}

If we let $a$ be the smallest sector angle and $r$ be the difference between consecutive sector angles, then we have the angles $a, a+r, a+2r, \cdots. a+11r$. Use the formula for the sum of an arithmetic sequence and set it equal to 360, the number of degrees in a circle.

 
 \begin{align*} \frac{a+a+11r}{2}\cdot 12 &= 360\\ 2a+11r &= 60\\ a &= \frac{60-11r}{2} \end{align*}


 All sector angles are integers so $r$ must be a multiple of 2. Plug in even integers for $r$ starting from 2 to minimize $a.$ We find this value to be 4 and the minimum value of $a$ to be $\frac{60-11(4)}{2} = \boxed{\textbf{(C)}\ 8}$

 



\textbf{Solution 3}

Starting with the smallest term, $a - 5x \cdots a, a + x \cdots a + 6x$ where $a$ is the sixth term and $x$ is the difference. The sum becomes $12a + 6x = 360$ since there are $360$ degrees in the central angle of the circle. The only condition left is that the smallest term in greater than zero. Therefore, $a - 5x > 0$.
 \[2a + x = 60\]
 \[x = 60 - 2a\]
 \[a - 5(60 - 2a) > 0\]
 \[11a > 300\]Since $a$ is an integer, it must be $28$, and therefore, $x$ is $4$. $a - 5x$ is $\boxed{\textbf{(C)}\ 8}$

 



\textbf{Solution 4 (I AM SPEED)}

We can see that $12 - 5 = 7$, this is less than $11$(obviously). This mean if we subtract the answer choices from $360$ (all angles add to $360$degrees because it is a cricle), the one that is divisible by 11 is the correct answer.

 
 \[360 - 5 = 355\]
 \[360 - 6 = 354\]
 \[360 - 8 = 352\]
 \[360 - 10 = 350\]
 \[360 - 12 = 348\]

 By using the $11$ divisibility rule, we can see that 352 is divisible by 11. Therefore the answer is $\boxed{\textbf{(C)}\ 8}$

 


	\item (\textit{AMC 10 10A 2012}) Externally tangent circles with centers at points $A$ and $B$ have radii of lengths $5$ and $3$, respectively. A line externally tangent to both circles intersects ray $AB$ at point $C$. What is $BC$?


 $\textbf{(A)}\ 4\qquad\textbf{(B)}\ 4.8\qquad\textbf{(C)}\ 10.2\qquad\textbf{(D)}\ 12\qquad\textbf{(E)}\ 14.4$


 
 \textbf{Answer: }D



\textbf{Solution}


                    \begin{figure}[H]
                        \centering
                        \includegraphics[width=0.5\textwidth]{Solution/AMC_10_10A_2012_11_Solution_Figure_1.png}
                    \end{figure}
                    

 Let $D$ and $E$ be the points of tangency on circles $A$ and $B$ with line $CD$. $AB=8$. Also, let $BC=x$. As $\angle ADC$ and $\angle BEC$ are right angles (a radius is perpendicular to a tangent line at the point of tangency) and both triangles share $\angle ACD$, $\triangle ADC \sim \triangle BEC$. From this we can get a proportion.

 $\frac{BC}{AC}=\frac{BE}{AD} \rightarrow \frac{x}{x+8}=\frac{3}{5} \rightarrow 5x=3x+24 \rightarrow x=\boxed{\textbf{(D)}\ 12}$

 


 


	\item (\textit{AMC 10 10A 2012}) A year is a leap year if and only if the year number is divisible by $400$ (such as $2000$) or is divisible by $4$ but not $100$ (such as $2012$).  The $200$th anniversary of the birth of novelist Charles Dickens was celebrated on February $7$, $2012$, a Tuesday.  On what day of the week was Dickens born?


 $\textbf{(A)}\ \text{Friday}\qquad\textbf{(B)}\ \text{Saturday}\qquad\textbf{(C)}\ \text{Sunday}\qquad\textbf{(D)}\ \text{Monday}\qquad\textbf{(E)}\ \text{Tuesday}$


 
 \textbf{Answer: }A



\textbf{Solution}

In this solution we refer to moving to the left as decreasing the year or date number and moving to the right as  increasing the year or date number. Every non-leap year we move to the right results in moving one day to the right because $365\equiv 1\pmod 7$.  Every leap year we move to the right results in moving $2$ days to the right since $366\equiv 2\pmod 7$. A leap year is usually every four years, so 200 years would have $\frac{200}{4}$ = $50$ leap years, but the problem says that 1900 does not count as a leap year.

 Therefore there would be 151 regular years and 49 leap years, so $1(151)+2(49)$ = $249$ days back.  Since $249 \equiv 4\ (\text{mod}\ 7)$, four days back from Tuesday would be $\boxed{\textbf{(A)}\ \text{Friday}}$.

 



\textbf{Solution 2}

Since they say that February $7$th, $2012$ is the $200$th anniversary of Charles Dickens birthday, that means that the birth of Charles Dickens is on February $7$th, $1812$. We then see that there is a leap year on $1812, 1816, ...., 2008$ but we must exclude $1900$ which equates to $49$ leap years. So, the amount of days we have to go back is $200(365) + 49$ days which in $\text{mod }7$ gives us 4. Thus, $4$ days back from Tuesday is $\boxed{\textbf{(A)}\ \text{Friday}}$.

 


 \textbf{Remark} The reason why there are 49 leap years:

 $\quad$ There are 200 $\div$ 4  years which are divisible by 4. There are 2 years which are divisible by 100 within these 200 years. There is 1 year which is divisible by 400. Therefore , there are $\frac{200}{4} -2 + 1$ leap years.

 


	\item (\textit{AMC 10 10A 2012}) An \textit{iterative average} of the numbers $1$, $2$, $3$, $4$, and $5$ is computed the following way. Arrange the five numbers in some order. Find the mean of the first two numbers, then find the mean of that with the third number, then the mean of that with the fourth number, and finally the mean of that with the fifth number. What is the difference between the largest and smallest possible values that can be obtained using this procedure?


 $\textbf{(A)}\ \frac{31}{16}\qquad\textbf{(B)}\ 2\qquad\textbf{(C)}\ \frac{17}{8}\qquad\textbf{(D)}\ 3\qquad\textbf{(E)}\ \frac{65}{16}$


 
 \textbf{Answer: }C



\textbf{Solution 1}

The iterative average of any 5 integers $a,b,c,d,e$ is defined as:

 
 \[\frac{\frac{\frac{\frac{a+b} 2+c} 2+d} 2+e} 2=\frac{a+b+2c+4d+8e}{16}\]

 
Plugging in $1,2,3,4,5$ for $a,b,c,d,e$, we see that in order to maximize the fraction,

 $a=1,b=2,c=3,d=4,e=5$,

 and in order to minimize the fraction,

 $a=5,b=4,c=3,d=2,e=1$.

 After plugging in these values and finding the positive difference of the two fractions, we arrive with $\frac{34}{16} \Rightarrow \frac{17}{8}$, which is our answer of $\boxed{\textbf{(C)}}$

 



\textbf{Solution 2}

The minimum and maximum can be achieved with the orders $5, 4, 3, 2, 1$ and $1, 2, 3, 4, 5$ respectively. We can see this because the iterative average is like a weighted average that gives more weight to later numbers. 

 $5,4,3,2,1 \Rightarrow \frac92,3,2,1 \Rightarrow \frac{15}{4},2,1 \Rightarrow \frac{23}{8},1 \Rightarrow \frac{31}{16}$

 $1,2,3,4,5 \Rightarrow \frac32,3,4,5 \Rightarrow \frac94,4,5 \Rightarrow \frac{25}{8},5 \Rightarrow \frac{65}{16}$

 The difference between the two is $\frac{65}{16}-\frac{31}{16}=\frac{34}{16}=\boxed{\textbf{(C)}\ \frac{17}{8}}$.

 


	\item (\textit{AMC 10 10A 2012}) Chubby makes nonstandard checkerboards that have $31$ squares on each side.  The checkerboards have a black square in every corner and alternate red and black squares along every row and column.  How many black squares are there on such a checkerboard?


 $\textbf{(A)}\ 480 \qquad\textbf{(B)}\ 481 \qquad\textbf{(C)}\ 482 \qquad\textbf{(D)}\ 483 \qquad\textbf{(E)}\ 484$


 
 \textbf{Answer: }B



\textbf{Solution 1}

There are 15 rows with 15 black tiles, and 16 rows with 16 black tiles, so the answer is $15^2+16^2 =225+256= \boxed{\textbf{(B)}\ 481}$

 Note: When solving $16^2$+$15^2$, you only need to calculate the units digit.

 



\textbf{Solution 2}

We build the $31 \times 31$ checkerboard starting with a board of $30 \times 30$ that is exactly half black.  There are $15 \cdot 30$ black tiles in this region.

 Add to this $30 \times 30$ checkerboard a $1 \times 30$ strip on the bottom that has $15$ black tiles.

 Add to this $31 \times 30$ checkerboard a $31 \times 1$ strip on the right that has $15 + 1$ black tiles.

 In total, there are $15 \cdot 30 + 15 + 15 + 1 = 481$ tiles, giving an answer of $\boxed{\textbf{(B)}\ 481}$

 



\textbf{Solution 3}

For every $2 \times 31$ strip, there are 31 black tiles and 31 white tiles. There are 15 such strips in total and another $1 \times 31$ strip having 16 black titles.

 In total, there are $31\times 15+16=481$ black tiles.

 



\textbf{Solution 4}

Drawing smaller scale sketches, we notice that the odd columns of an $n \times n$ (where $n$ is odd) board have $\left \lceil \frac{n}{2} \right \rceil$ black tiles, while the even columns have $\left \lfloor \frac{n}{2} \right \rfloor$ black tiles. In our case, we have a $31 \times 31$ board. We have $16$ odd columns, and $15$ even columns, so the number of black tiles in total is $16^2 + 15^2 = \boxed{\textbf{(B)}\ 481}$.

 



\textbf{Solution 5 (intuitive, fun)}

Consider a “snake” that begins at the leftmost vertex and travels down and up. In the first column, there are 16 black squares, and in the second, there are 15. This process repeats for the remaining rows, until the last, in which we end with 16. So, the answer is 
 \[16^2 + 15^2 = \boxed{481}\].

 


	\item (\textit{AMC 10 10A 2012}) Three unit squares and two line segments connecting two pairs of vertices are shown. What is the area of $\triangle ABC$?


 
                    \begin{figure}[H]
                        \centering
                        \includegraphics[width=0.5\textwidth]{Problem/AMC_10_10A_2012_15_Problem_Figure_1.png}
                    \end{figure}
                    

 $\textbf{(A)}\ \frac16 \qquad\textbf{(B)}\ \frac15 \qquad\textbf{(C)}\ \frac29 \qquad\textbf{(D)}\ \frac13 \qquad\textbf{(E)}\ \frac{\sqrt{2}}{4}$


 
 \textbf{Answer: }B



\textbf{Solution 1}


                    \begin{figure}[H]
                        \centering
                        \includegraphics[width=0.5\textwidth]{Solution/AMC_10_10A_2012_15_Solution_Figure_1.png}
                    \end{figure}
                    

 $AC$ intersects $BC$ at a right angle, (this can be proved by noticing that the slopes of the two lines are negative reciprocals of each other) so $\triangle ABC \sim \triangle BED$. The hypotenuse of right triangle $BED$ is $\sqrt{1^2+2^2}=\sqrt{5}$.

 
 \[\frac{AC}{BC}=\frac{BD}{ED} \Rightarrow \frac{AC}{BC} = \frac21 \Rightarrow AC=2BC\]

 
 \[\frac{AC}{AB}=\frac{BD}{BE} \Rightarrow \frac{AC}{1}=\frac{2}{\sqrt{5}} \Rightarrow AC=\frac{2}{\sqrt{5}}\]

 Since 

 


	\item (\textit{AMC 10 10A 2012}) Three runners start running simultaneously from the same point on a $500$-meter circular track. They each run clockwise around the course maintaining constant speeds of $4.4$, $4.8$, and $5.0$ meters per second. The runners stop once they are all together again somewhere on the circular course. How many seconds do the runners run?


 $\textbf{(A)}\ 1,000\qquad\textbf{(B)}\ 1,250\qquad\textbf{(C)}\ 2,500\qquad\textbf{(D)}\ 5,000\qquad\textbf{(E)}\ 10,000$


 
 \textbf{Answer: }C



\textbf{Solution 1}

First consider the first two runners. The faster runner will lap the slower runner exactly once, or run 500 meters farther. Let $x$ be the time these runners run in seconds.

 
 \[4.8x-4.4x=500 \Rightarrow x=1250\]

 Because $4.4(1250)=5500$ is a multiple of 500, it turns out they just meet back at the start line.

 Now we must find a time that is a multiple of $1250$ and results in the 5.0 m/s runner to end up on the start line. Every $1250$ seconds, that fastest runner goes $5.0(1250)=6250$ meters. In $2(1250)=2500$ seconds, he goes $5.0(2500)=12500$ meters. Therefore the runners run $\boxed{\textbf{(C)}\ 2,500}$ seconds.

 



\textbf{Solution 2}

Working backwards from the answers starting with the smallest answer, if they had run $1000$ seconds, they would have run $4400, 4800, 5000$ meters, respectively.  The first two runners have a difference of $400$ meters, which is not a multiple of $500$ (one lap), so they are not in the same place.

 If they had run $1250$ seconds, the runners would have run $5500, 6000, 6250$ meters, respectively.  The last two runners have a difference of $250$ meters, which is not a multiple of $500$.

 If they had run $2500$ seconds, the runners would have run $11000, 12000, 12500$ meters, respectively.  The distance separating each pair of runners is a multiple of $500$, so the answer is $\boxed{\textbf{(C)}\ 2,500}$ seconds.

 



\textbf{Solution 3}

Let $t$ be the time run in seconds, then the difference in meters run between the three runners is $0.2t, 0.4t, 0.6t$. For them to be at the same location all of them need to be multiples of 500. It is now easy to see that $0.2t=500, 0.4t=1000, 0.6t=1500$, so $t=\boxed{\textbf{(C)}\ 2,500}$.

 



\textbf{Solution 4}

After $t$ seconds, respectively the runners would've ran $4.4t, 4.8t,$ and $5t$ meters. Their current positions on the track are these values $\pmod{500}$. We're trying to find the value of $t$ such that 
 \[4.4t \equiv 4.8t \equiv 5t \pmod{500}\] Subtracting $4.4t$ on all sides, we get 
 \[0 \equiv 0.4t \equiv 0.6t \pmod{500}\] Now, we must find a value for $t$ such that both $0.6t$ and $0.4t$ are simultaneously multiples of $500$.

 Plugging in $500$ for $0.4t$ we get $t=1250$, but this does not work for $0.6t$ ($750$ isn't a multiple of $500$). Plugging in $0.4t=1000$, we get $t=2500$, and this does work for $0.6t$. 

 Therefore, $t=2500$ and the answer is $\boxed{\textbf{(C) } 2500}$.

 


 
                \begin{itemize}
                    \item Note: Modular Arithmetic works only for integral values, so my usage of decimals is technically incorrect but the intuition leads to the right answer
                \end{itemize}


                



\textbf{Solution 5}

Similar to the solution above, but is much quicker and does not involve trial and error. This uses decimal mod arithmetic, which can be justified by intuition...After $t$ seconds, respectively the runners would've ran $4.4t, 4.8t,$ and $5t$ meters. These three values are congruent $\pmod{500}$, so 
 \[4.4t \equiv 4.8t \equiv 5t \pmod{500}\]. Subtract $4.4t$ from all three sides to get $0, 0.4t,$ and $0.6t$ are congruent. Now all we need to find is a value of $t$ for which $0.4t$ and $0.6t$ are congruent $\pmod{500}$. Subtract $0.4t$ from both sides to get $0.2t$ and $0$ are congruent mod $500$, or that $0.2t=\dfrac{t}{5}$ is a multiple of $500$. Let $t=500k$, so we want $100k$ to be a multiple of $500$, or $k$ to be a multiple of $5$. Therefore, the smallest value of $t$ is when $k=5$, and when $t=500k=500(5)=2500 \space \boxed{(\text{C})}$

 



\textbf{Solution 6 (guessing)}

First, we determine the LCM of 44, 48, and 50, which is 600(44) = 26,400. Estimating this (ie. division by 10) gives 2640, and from this, we estimate that the time is 2500 seconds. $\boxed{(\text{C})}$ ~elpianista227

 


	\item (\textit{AMC 10 10A 2012}) Let $a$ and $b$ be relatively prime positive integers with $a>b>0$ and $\dfrac{a^3-b^3}{(a-b)^3} = \dfrac{73}{3}$. What is $a-b$?


 $\textbf{(A)}\ 1\qquad\textbf{(B)}\ 2\qquad\textbf{(C)}\ 3\qquad\textbf{(D)}\ 4\qquad\textbf{(E)}\ 5$


 
 \textbf{Answer: }C



\textbf{Solution 1 (Quick Insight)}

Since $a$ and $b$ are relatively prime, $a^3-b^3$ and $(a-b)^3$ are both integers as well. Then, for the given fraction to simplify to $\frac{73}{3}$, the denominator $(a-b)^3$ must be a multiple of $3.$ Thus, $a-b$ is a multiple of $3$. Looking at the answer choices, the only multiple of $3$ is $\boxed{\textbf{(C)}\ 3}$.

 



\textbf{Solution 2}

Using difference of cubes in the numerator and cancelling out one $(a-b)$ in the numerator and denominator gives $\frac{a^2 + ab + b^2}{a^2 - 2ab + b^2} = \frac{73}{3}$.

 Set $x = a^2 + b^2$, and $y = ab$.  Then $\frac{x + y}{x - 2y} = \frac{73}{3}$.  Cross multiplying gives 

 


	\item (\textit{AMC 10 10A 2012}) The closed curve in the figure is made up of $9$ congruent circular arcs each of length $\frac{2\pi}{3}$, where each of the centers of the corresponding circles is among the vertices of a regular hexagon of side $2$. What is the area enclosed by the curve?


 
                    \begin{figure}[H]
                        \centering
                        \includegraphics[width=0.5\textwidth]{Problem/AMC_10_10A_2012_18_Problem_Figure_1.png}
                    \end{figure}
                    


 $\textbf{(A)}\ 2\pi+6\qquad\textbf{(B)}\ 2\pi+4\sqrt{3}\qquad\textbf{(C)}\ 3\pi+4\qquad\textbf{(D)}\ 2\pi+3\sqrt{3}+2\qquad\textbf{(E)}\ \pi+6\sqrt{3}$


 
 \textbf{Answer: }E



\textbf{Solution 1}


                    \begin{figure}[H]
                        \centering
                        \includegraphics[width=0.5\textwidth]{Solution/AMC_10_10A_2012_18_Solution_Figure_1.png}
                    \end{figure}
                    

 We can draw the hexagon between the centers of the circles, and compute its area. The hexagon is made of $6$ equilateral triangles each with length $2$, so the area is:
 \[\frac{\sqrt{3}}{4} \cdot 2^2 \cdot 6=6 \sqrt{3}.\]Then, we add the areas of the three sectors outside the hexagon:
 \[\frac 23 \pi \cdot 3=2\pi.\]We now subtract the areas of the three sectors inside the hexagon but outside the figure (which is $\pi$) to get the area enclosed in the curved figure:
 \[6 \sqrt{3}+2\pi-\pi=\pi+6\sqrt{3}.\]Hence, our answer is $\boxed{\textbf{(E)}\ \pi+6\sqrt{3}},$ and we are done.
 \[\](Minor edits, made by dbnl.)

 



\textbf{Solution 2 (Looking at the answer choices)}

After forming the hexagon using the sectors outside the hexagon, we see three sectors left. Each sector has an area of $\frac{\pi}{3},$ so the three combined make $\pi.$ Since the side length of the hexagon is $2,$ its area doesn't have $\pi$ in it, so we know that the final answer will be $\pi + \text{(area of the hexagon)}.$ Looking at the answer choices, the only answer with only one $\pi$ is $\boxed{\textbf{(E)}}.$

 Also, notice the problem consists of adding and subtracting arcs from a hexagon with an area of $6\sqrt{3}$. Since all of the arcs have $\pi$ in them, they will not affect the area of the hexagon, which is $6\sqrt{3}$, in the final answer. Thus, the only possible solution is $\boxed{E}$. ~Extremelysupercooldude

 



\textbf{Solution 3 (Areas)}

As you can see, this diagram looks like a fidget spinner ;). Fidget spinners aside, we need to add stuff to our diagram to make it look easier. In the directions, they were talking about the centers of each arc create a hexagon, so let's add that to our diagram. 

 
                    \begin{figure}[H]
                        \centering
                        \includegraphics[width=0.5\textwidth]{Solution/AMC_10_10A_2012_18_Solution_Figure_2.png}
                    \end{figure}
                    

 The side length of the hexagon is 2 and if we plug it in to the area of a regular hexagon formula $\frac{3\sqrt 3}{2}s^2$ we get $6\sqrt 3$.

 Note that the interior angles of a regular hexagon is 120 because of the formula $\frac{180(n-2)}{n}$ where n is the number of sides. Knowing that, each sector is $\frac{1}{3}$ of a circle. This would mean the three sectors inside the hexagon altogether equal a full circle. Knowing that the radius is 1, we can use the area of a circle $\pi r^2$ and subtract it to $6\sqrt 3$. Thus we get the total area of $6\sqrt 3 - \pi$.

 
Notice that we have three sectors exterior to the hexagon. Realize that the central angles of a circle always sum up to 360. Since we know one of the central angles is equal to 120, then we subtract 360 to 120 which gives us 240. Knowing that, each sector is $\frac{2}{3}$ of a circle and since there is 3 of them, $\frac{2}{3}*3=2$ circles. To find the area of those circles, we have to use $\pi r^2$ again, but since there are 2 circles, then it is $2\pi r^2$, which gives us $2\pi$. 

 Now we have enough information to find the total area, 

 $(6\sqrt 3 -\pi+2\pi)=\textbf{(E)}\ \pi+6\sqrt{3}$

 



\textbf{Solution 4 (Quick Equilateral Triangles)}

First, realize that if you put 3 points around the tips of the fidget spinner shape, you get a big equilateral triangle, like this:
                    \begin{figure}[H]
                        \centering
                        \includegraphics[width=0.5\textwidth]{Solution/AMC_10_10A_2012_18_Solution_Figure_3.png}
                    \end{figure}
                    

 This cleanly gives the 9 arcs, and some can be cancelled out. 3 of these arcs make a circle, and each of the 3 arcs are $\frac{2\pi}{3}$. This means that when 3 of the small arcs are combined, it gives the circumference, which is $\frac{2\pi}{3}\cdot3=2\pi$. This means the radius of the small circles are each $\frac{2}{2}=1$.
                    \begin{figure}[H]
                        \centering
                        \includegraphics[width=0.5\textwidth]{Solution/AMC_10_10A_2012_18_Solution_Figure_4.png}
                    \end{figure}
                    

 It can be realized that the large equilateral triangle can be divided into 9 smaller equilateral trianlges, and they are inscribed in the small circles, as shown below.
                    \begin{figure}[H]
                        \centering
                        \includegraphics[width=0.5\textwidth]{Solution/AMC_10_10A_2012_18_Solution_Figure_5.png}
                    \end{figure}
                    

 As shown, each of the small equilateral triangles are inscribed in the small circles. Since we know the radius of the circle, using the formula for equilateral triangle inscribed in a circle, the area of the small equilateral triangle is $\frac{3}{4}\sqrt{3}\cdot1^2$. To solve the actual problem, you need to take the area of the equilateral triangle, and add it to the area of 1 small circle, subtracted from the area of 1 small equilateral triangle (all the other small arc figures cancel out, as shown in the past diagrams). The area of the large equilateral triangle is $9$ $\cdot$ $\frac{3}{4}\sqrt{3}\cdot1^2$ (Latter is area of each small equilateral triangle.) This equals $\frac{27}{4}\sqrt{3}\cdot1^2$. We know the area of 1 small circle is $\pi$ (we found radius earlier), so subtracting the area of the small equilateral triangle gives $\pi$ $-$ $\frac{3}{4}\sqrt{3}\cdot1^2$. Then adding this back to the large equilateral triangle (we already subtracted earlier), gives $\frac{27}{4}\sqrt{3}\cdot1^2+\pi-\frac{3}{4}\sqrt{3}\cdot1^2$. This gives $\frac{24}{4}\sqrt{3}+\pi$, and when simplified, $6\sqrt{3}+\pi$ $=$ $\boxed{\textbf{(E)}\ \pi+6\sqrt{3}}$.

 Small Alternate Solution - Since it is seen that there is $\pi$ when we simplify the expressions above, the only answer is $\boxed{\textbf{(E)}\ \pi+6\sqrt{3}}$.

 


	\item (\textit{AMC 10 10A 2012}) Paula the painter and her two helpers each paint at constant, but different, rates. They always start at 8:00 AM, and all three always take the same amount of time to eat lunch. On Monday the three of them painted 50\% of a house, quitting at 4:00 PM. On Tuesday, when Paula wasn't there, the two helpers painted only 24\% of the house and quit at 2:12 PM. On Wednesday Paula worked by herself and finished the house by working until 7:12 PM. How long, in minutes, was each day's lunch break?


 $\textbf{(A)}\ 30\qquad\textbf{(B)}\ 36\qquad\textbf{(C)}\ 42\qquad\textbf{(D)}\ 48\qquad\textbf{(E)}\ 60$


 
 \textbf{Answer: }D



\textbf{Solution 1}

Let Paula work at a rate of $p$, the two helpers work at a combined rate of $h$, and the time it takes to eat lunch be $L$, where $p$ and $h$ are in house/hours and L is in hours. Then the labor on Monday, Tuesday, and Wednesday can be represented by the three following equations:

 
 \[(8-L)(p+h)=50\]

 
 \[(6.2-L)h=24\]

 
 \[(11.2-L)p=26\]

 With three equations and three variables, we need to find the value of $L$.Adding the second and third equations together gives us $6.2h+11.2p-L(p+h)=50$. Subtracting the first equation from this new one gives us $-1.8h+3.2p=0$, so we get $h=\frac{16}{9}p$. Plugging into the second equation:

 
 \[(6.2-L)\frac{16}{9}p=24\]
 \[(6.2-L)p=\frac{27}{2}\]

 We can then subtract this from the third equation:

 
 \[5p=26-\frac{27}{2}\]
 \[p=\frac{5}{2}\]Plugging $p$ into our third equation gives: 
 \[L=\frac{4}{5}\]

 Converting $L$ from hours to minutes gives us $L=48$ minutes, which is $\boxed{\textbf{(D)}\ 48}$.

 



\textbf{Solution 2 (Modular Arithmetic)}

Because Paula worked from 
 \[8:00 \text{A.M.}\] to 
 \[7:12 \text{P.M.}\], she worked for 11 hours and 12 minutes = 672 minutes. Since there is $100-50-24=26$\% of the house left, we get the equation $26a=672$. Because $672$ is $22$ mod $26$, looking at our answer choices, the only answer that is $22$ $\text{mod}$ $26$ is $48$. So the answer is $\boxed{\textbf{(D)}\ 48}$.

 



\textbf{Solution 3 (GCD)}

We factor out the equations to be $(8-n)\left(\frac{1}{p}+\frac{1}{h}\right)=\frac{1}{2},\left(\frac{31}{5}-n\right)\left(\frac{1}{h}\right)=\frac{6}{25} \text{ and } \left(\frac{56}{5}-n\right)\left(\frac{1}{p}=\frac{13}{50}\right)$, where n is the number of hours for the break, p is the time Paula requires, and h is the time her helpers require. We find that when we select $\mathbf{D}$, we have them being $10.4$ and $5.4$, which correspond to being multiples of 13 and 6. Checking, we find that this satisfies the first equation, so multiplying $0.8\cdot 60= \boxed{48.}$

 


	\item (\textit{AMC 10 10A 2012}) A $3 \times 3$ square is partitioned into $9$ unit squares. Each unit square is painted either white or black with each color being equally likely, chosen independently and at random. The square is then rotated $90\,^{\circ}$ clockwise about its center, and every white square in a position formerly occupied by a black square is painted black. The colors of all other squares are left unchanged. What is the probability the grid is now entirely black?


 $\textbf{(A)}\ \frac{49}{512}\qquad\textbf{(B)}\ \frac{7}{64}\qquad\textbf{(C)}\ \frac{121}{1024}\qquad\textbf{(D)}\ \frac{81}{512}\qquad\textbf{(E)}\ \frac{9}{32}$


 
 \textbf{Answer: }A



\textbf{Solution 1}

First, look for invariants. The center, unaffected by rotation, must be black. So automatically, the chance is less than $\frac{1}{2}.$ Note that a $90^{\circ}$ rotation requires that black squares be across from each other across a vertical or horizontal axis. 

 
As such, $2$ squares directly across from each other must be black in the $4$ edge squares. Since there are $2$ configurations for this to be possible (top and bottom, right and left), this is a chance of
 \[\left (\frac{1}{2} \cdot \frac{1}{2} \right )+ \left (\frac{1}{2} \cdot \frac{1}{2} \right )=\frac{1}{2}\] 

 
However, by \href{https://artofproblemsolving.com/wiki/index.php?title=PIE}{PIE}, subtract the chance all 4 are black and both configurations are met: $\frac{1}{2}- \left (\frac{1}{2} \cdot \frac{1}{2} \right ) \cdot \left (\frac{1}{2} \cdot \frac{1}{2} \right )=\frac{7}{16}$. Through symmetrical reasoning, the corners also have a $\frac{7}{16}$ chance of having a configuration that yields all black corners. 

 Then, the chance that all squares black is the intersection of all these probabilities: $\frac{1}{2} \left (\frac{7}{16} \right ) \left (\frac{7}{16} \right ) = \boxed{\textbf{(A)}\ \frac{49}{512}}$ 

 


 Also, if you have little to no time and are guessing, notice there are a total of $2^9 = 512$ ways to permutate the colors on the square (Each square can be white or black, so there are 2 possibilities for each of the 9 squares). Thus, the answer must be in some form of $\frac{\text{the number of good cases}}{\text{the total number of cases}}$, so E is not possible. Also, since the number of good cases must be an integer, C is not possible. From there, your chances of guessing the right answer are slightly higher. ~Extremelysupercooldude

 



\textbf{Solution 2}

First, there is only one way for the middle square to be black because it is not affected by the rotation. Then we can consider the corners and edges separately. Let's first just consider the number of ways we can color the corners. There is $1$ case with all-black squares. There are four cases with one white square and all $4$ work. There are six cases with two white squares, but only the $2$ with the white squares diagonal from each other works. There are no cases with three white squares or four white squares. Then the total number of ways to color the corners is $1+4+2=7$. In essence, the edges work the same way, so there are also $7$ ways to color them. The number of ways to fit the conditions over the number of ways to color the squares is

 
 \[\frac{7\times7}{2^9}=\boxed{\textbf{(A)}\ \frac{49}{512}}\]

 



\textbf{Solution 3}

We proceed by casework.Note that the middle square must be black because when rotated 90 degrees, it must keep its position. Now we have to deal with the following cases: 
Case 1: 0 white squares.There is exactly $1$ way to color the grid this way. 
Case 2: 1 white square.Note that the white square can be anywhere on the grid except for the middle square because when rotating 90 degrees it can never land on itself. Thus, there are $8$ cases.
Case 3: 2 white squares. We have $\binom{8}{2}=28$ ways to color two white squares without restrictions (the middle square must be black, giving us 8 squares to choose from). However, we must subtract the ways where two white squares differ by a rotation of 90 degrees about the middle of the square. There are a total of 8 cases we must subtract (these are not too hard to see). Thus, there are $20$ ways from this.
Case 4: 3 white squares.Since we can not change the middle square, there are $\binom{8}{3}=56$ ways to color this. However, we must consider the cases where at least two squares differ by a rotation of 90 degrees. We can count this with PIE: by the Principle of Exclusion, the number of cases we want to exclude is the number of cases where 2 squares differ by a rotation of 90 degrees and minus when there are 3 squares such that two of them differ by rotation of 90 and 1 of them differs by rotation of 180, because of the overcount from our first case. From case 2, there are 8 causes such that two squares differ by a rotation of 90. There are also 6 other places we can place the third square (it can't be the middle of the two that we already colored), for a total of 48 ways. We have to subtract the second case. Note that there are 8 ways in which we can arrange two squares differing by 180 degrees. Out of these, each one has two ways to put another square such that two differ by 90 degrees and 1 pair differs by 180. However, this is overcounted by a factor of 2, so there are act/2=8 ways. Thus, we have $56-(48-8)=16$ ways in this case.
Case 5: 4 white squares.Note that two of them have to be on one of the 4 corner squares, and two of them have to be on one of the 4 edge squares. Each solution yields two combinations, for a total of 2 * 2 = 4.Adding up our cases yields 1+8+20+16+4=49 ways. There are 512 ways to color the square without restrictions. Thus, the answer is $\boxed{\textbf{(A)}\ 49/512}$.

 



\textbf{Solution 4(Similar to Solution 2)}

Number the squares in the original grid $1,2,3$ for the first row then $4,5,6$ for the second and lastly $7,8,9$ for the last. Essentially the grid should look like this:

 1 2 3

 4 5 6

 7 8 9

 Next, consider the numbering of the grid after a $90$ degree clockwise rotation about the center. This is pretty easy to visualize. The grid should end up looking like:

 7 4 1

 8 5 2

 9 6 3

 I define an orbit to be a cycle of the squares affected by the previous square's color. For instance, a corner orbit would be 1 - 7 - 9 - 3 - 1 because after the rotation, square $7$ is getting affected by square $1$'s position and we can reason for the others. Next, an edge orbit would be 2 - 4 - 8 - 6 - 2. We will first find the total number of colored orbits for a given corner orbit. The number of cases we will get for a corner orbit would be the same as the number of cases we will get for an edge orbit. We will multiply these cases together, and divide by the total number of configurations which is $2^{9} = 512$. A corner orbit is 1 - 7 - 9 - 3 - 1. If square $1$ was originally colored black, then square $7$ has absolutely no restriction at all because it's transformed position is exactly where square $1$ was, which was colored black. If square $7$ was left colored white, then we can easily see that square $9$ then must be colored black. Finally, square $3$ has absolutely no restrictions at all. Let's assume it was colored white originally. Then one possible corner orbit would be B - W - B - B - B. In fact, in any given orbit, if two entries in an orbit are colored white, then the orbit wouldn't exist. This means that say we color both $7$ and $9$ white in a corner orbit. Then, the orbit wouldn't exist because the new grid would all be covered black there would still be some white squares. Bearing this in mind and checking if our hypothesis stands true, we can quickly list down the rest of the corner orbit cases:

 B - W - B - B - B

 B - W - B - W - B

 B - B - W - B - B

 B - B - B - W - B

 B - B - B - B - B

 W - B - W - B - W

 W - B - B - B - B

 As one can astutely notice, no "W" entries are sitting right next to each other! All of these orbits would actually work. Thus, we have a total of $7$ corner orbits. As we said before, because of the symmetry of a square, the edge orbits would function absolutely the same way as the corner orbits and thus we would also have $7$ edge orbits. To end with a full on black grid, we must combine a corner orbit with an edge orbit for a total of $7 \cdot 7 = 49$ cases. Thus the probability is $\frac{49}{512}$ which results to $\boxed{A}$. 

 Note: The $5$ square in the middle just stays the same even after the rotation so that is why it isn't included in any of the orbits. 

 


	\item (\textit{AMC 10 10A 2012}) Let points $A = (0 ,0 ,0)$, $B = (1, 0, 0)$, $C = (0, 2, 0)$, and $D = (0, 0, 3)$. Points $E$, $F$, $G$, and $H$ are midpoints of line segments $\overline{BD},\text{ }  \overline{AB}, \text{ } \overline {AC},$ and $\overline{DC}$ respectively. What is the area of $EFGH$?


 $\textbf{(A)}\ \sqrt{2}\qquad\textbf{(B)}\ \frac{2\sqrt{5}}{3}\qquad\textbf{(C)}\ \frac{3\sqrt{5}}{4}\qquad\textbf{(D)}\ \sqrt{3}\qquad\textbf{(E)}\ \frac{2\sqrt{7}}{3}$


 
 \textbf{Answer: }C



\textbf{Solution 1}

Consider a tetrahedron with vertices at $A,B,C,D$ in the $xyz$-space. The length of $EF$ is just one-half of $AD$ because it is the midsegment of $\triangle ABD.$ The same concept applies to the other side lengths. $AD=3$ and $BC=\sqrt{1^2+2^2}=\sqrt{5}$. Then $EF=HG=\frac32$ and $EH=FG=\frac{\sqrt{5}}{2}$. The line segments lie on perpendicular planes so quadrilateral $EFGH$ is a rectangle. The area is

 
 \[EF \cdot FG = \frac 32 \cdot \frac{\sqrt{5}}{2} = \frac{3\sqrt 5} 4\implies \boxed{\textbf C}.\]

 
                    \begin{figure}[H]
                        \centering
                        \includegraphics[width=0.5\textwidth]{Solution/AMC_10_10A_2012_21_Solution_Figure_1.png}
                    \end{figure}
                    

 



\textbf{Solution 2}

Computing the points of $EFGH$ gives $E(0.5, 0, 1.5), F(0.5, 0, 0), G(0,1,0), H(0,1,1.5)$.  The vector $EF$ is $(0,0,-1.5)$, while the vector $HG$ is also $(0,0,-1.5)$, meaning the two sides $EF$ and $GH$ are parallel.  Similarly, the vector $FG$ is $(-0.5, 1, 0)$, while the vector $EH$ is also $(-0.5, 1, 0)$.  Again, these are equal in both magnitude and direction, so $FG$ and $EH$ are parallel.  Thus, figure $EFGH$ is a parallelogram.

 Computation of vectors $EF$ and $HG$ is sufficient evidence that the figure is a parallelogram, since the vectors are not only point in the same direction, but are of the same magnitude, but the other vector $FG$ is needed to find the angle between the sides.

 Taking the dot product of vector $EF$ and vector $FG$ gives $0 \cdot -0.5 + 0 \cdot 1 + -1.5 \cdot 0 = 0$, which means the two vectors are perpendicular.  (Alternately, as above, note that vector $EF$ goes directly down on the z-axis, while vector $FG$ has no z-component and lie completely in the xy plane.)  Thus, the figure is a parallelogram with a right angle, which makes it a rectangle.  With the distance formula in three dimensions, we find that $EF = \frac{3}{2}$ and $FG = \frac{\sqrt{5}}{2}$, giving an area of $\frac32 \cdot \frac{\sqrt{5}}{2} = \boxed{\textbf{(C)}\ \frac{3\sqrt{5}}{4}}$

 


	\item (\textit{AMC 10 10A 2012}) The sum of the first $m$ positive odd integers is $212$ more than the sum of the first $n$ positive even integers. What is the sum of all possible values of $n$?


 $\textbf{(A)}\ 255\qquad\textbf{(B)}\ 256\qquad\textbf{(C)}\ 257\qquad\textbf{(D)}\ 258\qquad\textbf{(E)}\ 259$


 
 \textbf{Answer: }A



\textbf{Solution 1}

The sum of the first $m$ odd integers is given by $m^2$. The sum of the first $n$ even integers is given by $n(n+1)$.

 Thus, $m^2 = n^2 + n + 212$. Since we want to solve for n, rearrange as a quadratic equation: $n^2 + n + (212 - m^2) = 0$.

 
Use the quadratic formula: $n = \frac{-1 + \sqrt{1 - 4(212 - m^2)}}{2}$. Since $n$ is clearly an integer, $1 - 4(212 - m^2) = 4m^2 - 847$ must be not only a perfect square, but also an odd perfect square for $n$ to be an integer.

 
Let $x = \sqrt{4m^2 - 847}$; note that this means $n = \frac{-1 + x}{2}$. It can be rewritten as $x^2 = 4m^2 - 847$, so $4m^2 - x^2 = 847$. Factoring the left side by using the difference of squares, we get $(2m + x)(2m - x) = 847 = 7\cdot11^2$.

 
Our goal is to find possible values for $x$, then use the equation above to find $n$. The difference between the factors is 

 


	\item (\textit{AMC 10 10A 2012}) Adam, Benin, Chiang, DeShawn, Esther, and Fiona have internet accounts.  Some, but not all, of them are internet friends with each other, and none of them has an internet friend outside this group.  Each of them has the same number of internet friends.  In how many different ways can this happen?


 $\textbf{(A)}\ 60\qquad\textbf{(B)}\ 170\qquad\textbf{(C)}\ 290\qquad\textbf{(D)}\ 320\qquad\textbf{(E)}\ 660$


 
 \textbf{Answer: }B



\textbf{Solution}

Note that if $n$ is the number of friends each person has, then $n$ can be any integer from $1$ to $4$, inclusive, truly.

 One person can have at most 4 friends since they cannot be all friends (stated in the problem).

 Also note that the cases of $n=1$ and $n=4$ are the same, since a map showing a solution for $n=1$ can correspond one-to-one with a map of a solution for $n=4$ by simply making every pair of friends non-friends and vice versa.  The same can be said of configurations with $n=2$ when compared to configurations of $n=3$.  Thus, we have two cases to examine, $n=1$ and $n=2$, and we count each of these combinations twice.

 (Note: If you aren’t familiar with one-to-one correspondences, think of it like this: the number of ways to choose 4 friends is equal to number of ways to exclude one friend from your friend group. Hence, since the number of ways to choose 1 friend is the same thing as choosing 1 to not be friends with, $n=1$ and $n=4$ have the same number of ways. Similarly, $n=2$ and $n=3$ have the same number of ways as well. ~peelybonehead)

 For $n=1$, if everyone has exactly one friend, that means there must be $3$ pairs of friends, with no other interconnections.  The first person has $5$ choices for a friend.  There are $4$ people left.  The next person has $3$ choices for a friend.  There are two people left, and these remaining two must be friends.  Thus, there are $15$ configurations with $n=1$.

 For $n=2$, there are two possibilities.  The group of $6$ can be split into two groups of $3$, with each group creating a friendship triangle.  The first person has $\binom{5}{2} = 10$ ways to pick two friends from the other five, while the other three are forced together.  Thus, there are $10$ triangular configurations.  

 However, the group can also form a friendship hexagon, with each person sitting on a vertex, and each side representing the two friends that person has.  The first person may be seated anywhere on the hexagon \href{https://artofproblemsolving.com/wiki/index.php?title=Without\_loss\_of\_generality}{without loss of generality}.  This person has $\binom{5}{2} = 10$ choices for the two friends on the adjoining vertices.  Each of the three remaining people can be seated "across" from one of the original three people, forming a different configuration.  Thus, there are $10 \cdot 3! = 60$ hexagonal configurations, and in total $70$ configurations for $n=2$.

 As stated before, $n=3$ has $70$ configurations, and $n=4$ has $15$ configurations.  This gives a total of $(70 + 15)\cdot 2 = 170$ configurations, which is option $\boxed{\textbf{(B)}\ 170}$.

 
I thought the answer was 60??? \href{https://artofproblemsolving.com/wiki/index.php?title=User:Maxisw}{MisW} (\href{https://artofproblemsolving.com/wiki/index.php?title=User\_talk:Maxisw}{talk}) Reply: 60 is just for the subcase with one large cycle.

 \textbf{Note} We can also calculate the triangular configurations by applying $\frac{\binom{6}{3}}{2} = \frac{20}{2}=10$ (Because choosing $A$,$B$ and $C$ is the same as choosing $D$,$E$ and $F$.

 
For the hexagonal configurations, we know that the total amount of combinations is $6!=720$. However, we must correct for our  overcounting because of rotation and reflection. We have that there are ${720 \over 6 \cdot 2} = 60$ because there $6$ rotations of the hexagon and 2 reflections (clockwise and counterclockwise) for each valid way.

 \textbf{Note 2} We can also calculate the hexagonal configurations by placing each person at a random vertex (6!) and dividing for rotations (by 6) and reflections (by 2): 

 \begin{verbatim}
												$\frac{6!}{12}=60$
\end{verbatim}

\textbf{Note 3} For the hexagonal configuration, it is essentially congruent to counting the number of paths that start and end with A. From $A$, we have $5$ options, then the next one has $4$, and so on, so we have: $5\cdot 4 \hdots 1 \cdot 1$, where the final $1$ is because the final point needs to reconnect to $A$. Thus, there are $5! = 120$ ways. Dividing by $2$ since we overcounted by accounting for both directions, we have $120/2 = \fbox{60}$

 -bigbrain123

 \textbf{Yet Another Note (4)} You could ALSO count the hexagonal arrangements by considering the rotation and then reflection. Each arrangement of the six friends maps to a sequence \_ \_ \_ \_ \_ \_, each person is friends with their two neighbors, and the rightmost and leftmost people are friends with each other. Notice that any configuration can be rotated to begin with A, so we have the sequence: A \_ \_ \_ \_ \_. Now there are $5!$ ways to fill these blanks. For example, ABCDEF. However, since A is friends with both B and F, this sequence can be 'flipped' to yield AFEDCB, which corresponds to the same set of friendships. Since each sequence can be flipped to yield another valid sequence, we have overcounted by a factor of two. Thus the final number of hexagonal configurations = $\frac{5!}{2} = \fbox{60}$

 


	\item (\textit{AMC 10 10A 2012}) Let $a$, $b$, and $c$ be positive integers with $a\ge$ $b\ge$ $c$ such that

 \begin{align*}a^2-b^2-c^2+ab&=2011\text{ and}\\ a^2+3b^2+3c^2-3ab-2ac-2bc&=-1997.\end{align*}



 What is $a$?


 $\textbf{(A)}\ 249\qquad\textbf{(B)}\ 250\qquad\textbf{(C)}\ 251\qquad\textbf{(D)}\ 252\qquad\textbf{(E)}\ 253$


 
 \textbf{Answer: }E



\textbf{Solution}

Add the two equations.

 $2a^2 + 2b^2 + 2c^2 - 2ab - 2ac - 2bc = 14$.

 Now, this can be rearranged and factored.

 $(a^2 - 2ab + b^2) + (a^2 - 2ac + c^2) + (b^2 - 2bc + c^2) = 14$

 $(a - b)^2 + (a - c)^2 + (b - c)^2 = 14$

 $a$, $b$, and $c$ are all integers, so the three terms on the left side of the equation must all be perfect squares.We see that the only is possibility is $14 = 9 + 4 + 1$. 

 $(a-c)^2 = 9 \Rightarrow a-c = 3$, since $a-c$ is the biggest difference. It is impossible to determine by inspection whether $a-b = 1$ or $2$, or whether $b-c = 1$ or $2$.

 We want to solve for $a$, so take the two cases and solve them each for an expression in terms of $a$. Our two cases are $(a, b, c) = (a, a-1, a-3)$ or $(a, a-2, a-3)$. Plug these values into one of the original equations to see if we can get an integer for $a$.

 $a^2 - (a-1)^2 - (a-3)^2 + a(a-1) = 2011$, after some algebra, simplifies to$7a = 2021$. $2021$ is not divisible by $7$, so $a$ is not an integer. 

 The other case gives $a^2 - (a-2)^2 - (a-3)^2 + a(a-2) = 2011$, which simplifies to $8a = 2024$. Thus, $a = 253$ and the answer is $\boxed{\textbf{(E)}\ 253}$.

 


	\item (\textit{AMC 10 10A 2012}) Real numbers $x$, $y$, and $z$ are chosen independently and at random from the interval $[0,n]$ for some positive integer $n$. The probability that no two of $x$, $y$, and $z$ are within 1 unit of each other is greater than $\frac {1}{2}$. What is the smallest possible value of $n$?


 $\textbf{(A)}\ 7\qquad\textbf{(B)}\ 8\qquad\textbf{(C)}\ 9\qquad\textbf{(D)}\ 10\qquad\textbf{(E)}\ 11$


 
 \textbf{Answer: }D



\textbf{Solution 1}

Since $x,y,z$ are all reals located in $[0, n]$, the number of choices for each one is continuous so we use geometric probability.

 WLOG(\href{https://artofproblemsolving.com/wiki/index.php?title=Without\_loss\_of\_generality}{Without loss of generality}), assume that $n\geq x \geq y \geq z \geq 0$. Then the set of points $(x,y,z)$ is a tetrahedron, or a triangular pyramid. The point $(x,y,z)$ distributes uniformly in this region. If this is not easy to understand, read Solution II.

 The altitude of the tetrahedron is $n$ and the base is an isosceles right triangle with a leg length $n$. The volume is $V_1=\dfrac{n^3}{6}$, as shown in the first figure in red.

 
                    \begin{figure}[H]
                        \centering
                        \includegraphics[width=0.5\textwidth]{Solution/AMC_10_10A_2012_25_Solution_Figure_1.png}
                    \end{figure}
                    

 
Now we will find the region with points satisfying  $|x-y|\geq1$, $|y-z|\geq1$, $|z-x|\geq1$.

 Since $n\geq x \geq y \geq z \geq 0$, we have $x-y\geq1$, $y-z\geq1$.

 The region of points $(x,y,z)$ satisfying the condition is shown in the second figure in black. It is a tetrahedron, too.

 
                    \begin{figure}[H]
                        \centering
                        \includegraphics[width=0.5\textwidth]{Solution/AMC_10_10A_2012_25_Solution_Figure_2.png}
                    \end{figure}
                    

 The volume of this region is $V_2=\dfrac{(n-2)^3}{6}$.

 So the probability is $p=\dfrac{V_2}{V_1}=\dfrac{(n-2)^3}{n^3}$.

 Substituting $n$ with the values in the choices, we find that when $n=10$, $p=\frac{512}{1000}>\frac{1}{2}$, when $n=9$, $p=\frac{343}{729}<\frac{1}{2}$. So $n\geq 10$.

 So the answer is $\boxed{\textbf{(D)}\ 10}$.

 



\textbf{Solution 2}

Because $x$, $y$, and $z$ are chosen independently and at random from the interval $[0,n]$, which means that  $x$, $y$, and $z$ distributes uniformly and independently in the interval $[0,n]$. So the point $(x, y, z)$ distributes uniformly in the cubic $0\leqslant x, y, z \leqslant n$, as shown in the figure below. The volume of this cubic is $V_0=n^3$.

 
                    \begin{figure}[H]
                        \centering
                        \includegraphics[width=0.5\textwidth]{Solution/AMC_10_10A_2012_25_Solution_Figure_3.png}
                    \end{figure}
                    

 As we want to find the probablity of the incident $A=\big\{ |x-y| \geq 1, |y-z| \geq1, |z-x| \geq 1 \big\}$, we should find the volume of the region of points such that $|x-y|\geq 1$, $|y-z|\geq 1$, $|z-x|\geq 1$ and $0\leq x, y, z \leq n$.

 Now we will find the region  $\big\{ (x,y,z)\ | \ 0\leq x, y, z \leq n, |x-y|\geq 1, |y-z|\geq 1, |z-x|\geq 1 \big\}$.

 The region can be generated by cutting off 3 slices corresponding to $|x-y|< 1$, $|y-z|< 1$, and $|z-x|< 1$, respectively, from the cubic.

 After cutting off a slice corresponding to $|x-y|< 1$, we get two triangular prisms, as shown in the figure.

 
                    \begin{figure}[H]
                        \centering
                        \includegraphics[width=0.5\textwidth]{Solution/AMC_10_10A_2012_25_Solution_Figure_4.png}
                    \end{figure}
                    

 In order to observe the object clearly, we rotate the object by the $z$ axis, as shown.

 We can draw the slice corresponding to $|y-z|< 1$ on the object.

 
                    \begin{figure}[H]
                        \centering
                        \includegraphics[width=0.5\textwidth]{Solution/AMC_10_10A_2012_25_Solution_Figure_5.png}
                    \end{figure}
                    

 After cutting off the slice corresponding to $|y-z|< 1$, we have 4 pieces left.

 
                    \begin{figure}[H]
                        \centering
                        \includegraphics[width=0.5\textwidth]{Solution/AMC_10_10A_2012_25_Solution_Figure_6.png}
                    \end{figure}
                    

 After cutting off the slice corresponding to $|z-x|< 1$, we have 6 congruent triangular prisms. 

 
                    \begin{figure}[H]
                        \centering
                        \includegraphics[width=0.5\textwidth]{Solution/AMC_10_10A_2012_25_Solution_Figure_7.png}
                    \end{figure}
                    

 Here we draw all the pictures in colors in order to explain the solution clearly. That does not mean that the students should do it in the examination. They can draw a figure with lines only, as shown below.

 
                    \begin{figure}[H]
                        \centering
                        \includegraphics[width=0.5\textwidth]{Solution/AMC_10_10A_2012_25_Solution_Figure_8.png}
                    \end{figure}
                    

 Every triangular pyramid has an altitude $n-2$ and a base of isoceless right triangle with leg length $n-2$, so the volume is $(n-2)^3/6$.Then the volume of the region $\big\{ (x,y,z)\ | \ 0\leqslant x, y, z \leqslant n, |x-y|\geqslant 1, |y-z|\geqslant 1, |z-x|\geqslant 1 \big\}$ is $V_A=6\times(n-2)^3/6$=$(n-2)^3$.

 So the probability of the incident $A$ is $P(A)=\dfrac{V_A}{V_0}$=$\dfrac{(n-2)^3}{n^3}$.

 Then we can get the answer the same way as Solution I.

 The answer is $\boxed{\textbf{(D)}\ 10}$.

 


 
If there is no choice for selection, we can also find the minimum value of the integer $n$ if we do not substitute $n$ by the possible values one by one.

 Let $P(A)>1/2$, i.e., $\dfrac{(n-2)^3}{n^3}>\dfrac{1}{2}$,  so $\dfrac{n-2}{n}>\dfrac{1}{\sqrt[^3\!]{2}}$, or $1-\dfrac{2}{n}>\dfrac{1}{\sqrt[^3\!]{2}}$, hence $n>\dfrac{2\sqrt[^3\!]{2}}{\sqrt[^3\!]{2}-1}$.

 Now we will estimate the value of $\dfrac{2\sqrt[^3\!]{2}}{\sqrt[^3\!]{2}-1}$ without a calculator.

 Since $a^3-1$=$(a-1)(a^2+a+1)$, so$\dfrac{2\sqrt[^3\!]{2}}{\sqrt[^3\!]{2}-1}$=$\dfrac{2\sqrt[^3\!]{2}\times\left( \sqrt[^3\!]{2}^2+\sqrt[^3\!]{2}+1\right)}{\left( \sqrt[^3\!]{2}-1\right)\left( \sqrt[^3\!]{2}^2+\sqrt[^3\!]{2}+1\right)}$=$\dfrac{2\times\left( 2+\sqrt[^3\!]{2}^2+\sqrt[^3\!]{2}\right)}{ \sqrt[^3\!]{2}^3-1}$=$2\times\left( 2+\sqrt[^3\!]{4}+\sqrt[^3\!]{2}\right)$.

 Now we would get the approximation of $\sqrt[^3\!]{4}$ and $\sqrt[^3\!]{2}$.

 In order to avoid compicated computation, we get the approximation with one decimal digit only.

 Estimation of $\sqrt[^3\!]{2}$.

 Since $1.5^3=2.25\times1.5>2$, so $1<\sqrt[^3\!]{2}<1.5$.

 The mean of 1 and 1.5 with one decimal digit is about 1.3 .

 As $1.3^3=1.69\times 1.3=2.197>2$, so $1<\sqrt[^3\!]{2}<1.3$.

 The mean of 1 and 1.3 with one decimal digit is about 1.2.

 As $1.2^3=1.44\times 1.2=1.728<2$, so $1.2<\sqrt[^3\!]{2}<1.3$.

 Estimation of $\sqrt[^3\!]{4}$.

 As $\sqrt[^3\!]{4}=\sqrt[^3\!]{2}^2$, so $1.2^2<\sqrt[^3\!]{4}<1.3^2$,then $1.24<\sqrt[^3\!]{4}<1.69$.

 As $1.5^3=2.25\times 1.5=3.375<4$, so $1.5<\sqrt[^3\!]{4}<1.69$.

 The mean of 1.5 and 1.69 with one decimal digit is about 1.6.

 As $1.6^3=(16/10)^3=(2^4/10)^3=2^{12}/10^3=4\times 2^10/10^3=4\times 1.024>4$, so $1.5<\sqrt[^3\!]{4}<1.6$.

 
Then $2\times(2+1.5+1.2)<2\times\left(2+\sqrt[^3\!]{4}+\sqrt[^3\!]{2}\right)<2\times(2+1.6+1.3)$, i.e., $9.4<2\times\left(2+\sqrt[^3\!]{4}+\sqrt[^3\!]{2}\right)<9.8$,

 As $n>2\times\left(2+\sqrt[^3\!]{4}+\sqrt[^3\!]{2}\right)$, So the minimal value of integer $n$ is 10.

 \textbf{Appendix} This solution is motivated by the suggestive formula $\frac{(n-2)^{3}}{n^{3}}$.

 The problem generalizes easily to $k$-dimensional real space $\mathbb{R}^{k}$. In the general $k$-dimensional case, we are asked to find the probability that a randomly chosen $k$-tuple $(x_{1},\dotsc,x_{k}) \in [0,n]^{k}$ satisfies $|x_{i} - x_{j}| > 1$ for all $i \ne j$. To avoid repetition, let us say that $(x_{1},\dotsc,x_{k})$ is \textit{spaced-out} if $|x_{i} - x_{j}| > 1$ for all $i \ne j$.

 Let $C$ be the $k$-dimensional hyper-cube of side length $n$: 
 \[C = [0,n]^{k} = \big\{(x_{1},\dotsc,x_{k}) \in \mathbb{R}^{k} \;:\; 0 \leqslant x_{i} \leqslant n \text{ for all }i \big\} \;.\] Then $C$ has volume $n^{k}$. Let $S$ be the set of spaced-out $k$-tuples $(x_{1},\dotsc,x_{k})$. The desired probability is Vol$(S)/n^{k}$. 

 The set of $k$-tuples $(x_{1},\dotsc,x_{k})$ such that there exist distinct indices $i, j$ such that $x_{i} = x_{j}$ has volume $0$, so we may restrict our attention to $k$-tuples such that $x_{i} \ne x_{j}$ for all $i \ne j$. 

 Further, the condition that $(x_{1},\dotsc,x_{k})$ is spaced-out is "invariant upon permuting the indices"; in other words, if $\sigma$ is a permutation of the set of indices $\{1,\dotsc,k\}$, then $(x_{1},\dotsc,x_{k})$ is spaced-out if and only if $(x_{\sigma(1)},\dotsc,x_{\sigma(k)})$ is spaced-out. Therefore, we may consider the set of spaced-out $k$-tuples $(x_{1},\dotsc,x_{k})$ which additionally satisfy $x_{1} < \dotsb < x_{k}$. Let us denote this set by $T$. This condition is equivalent to 
 \[0 \leqslant x_{1} < x_{2} - 1 < \dotsb < x_{i}-(i-1) < \dotsb < x_{k}-(k-1) \leqslant n-(k-1) \;.\]Let us choose new variables $y_{i} = x_{i} - (i-1)$ for $i = 1,\dotsc,k$. This change of variables is just a translation of each $(x_{1},\dotsc,x_{k})$ by the vector $(0,1,\dots,k-1)$; in the above solutions, it corresponds to taking the 6 tetrahedrons and gluing them together to form a cube. 

 We now compute the volume of the set of $(y_{1},\dotsc,y_{k}) \in [0,n-(k-1)]^{k}$ which satisfy $y_{1} < \dotsb < y_{k}$. As above, we can disregard any $(y_{1},\dotsc,y_{k})$ such that $y_{i} = y_{j}$ for some $i \ne j$. Given any $(y_{1},\dotsc,y_{k})$ such that $y_{i} \ne y_{j}$ for all $i \ne j$, there exists exactly one permutation $\sigma$ of the indices $\{1,\dotsc,k\}$ such that $y_{\sigma(1)} < \dotsb < y_{\sigma(k)}$. Since there are $k!$ permutations of $\{1,\dotsc,k\}$, the desired volume is equal to $\frac{1}{k!}$ times the volume of the $k$-dimensional hyper-cube of side length $n-(k-1)$, which is $\frac{1}{k!}(n-(k-1))^{k}$. Hence $T$ has volume $\frac{1}{k!}(n-(k-1))^{k}$ as well and $S$ has volume $(n-(k-1))^{k}$. Hence the desired probability is $\frac{(n-(k-1))^{k}}{n^{k}}$.

 



\textbf{Solution 3}

If $x$, $y$, and $z$ are separated by at least one, then by subtracting the minimum space between the three variables (which is $2$), $x$, $y$, and $z$ can be chosen randomly in the interval $[0,n-2]$. The probability is hence $\dfrac{(n-2)^3}{n^3} > \dfrac{1}{2}.$ $\boxed{\textbf{(D)} 10}$ is the minimum value in the answer choices that satisfies this inequality.

 



\textbf{Solution 4}

Imagine Points $x$, $y$, $z$ as the "starting points" of three "blocks" of real numbers that have length $1$. We are just trying to find the probability that those three "blocks" do not overlap. To do this we can set each unit of $1$ into $\mu$ equal little increments, and take the limit of the probability as $\mu$ approaches $\infty$. (This is because there are indefinitely many real numbers in a given interval.) We see that the total number of arrangements for the three blocks of $1$, without considering the rule that no two blocks shall overlap is $(n\mu + 1)^3$/3!. We see that the number of ways to arrange the three blocks such that no three of them are overlapping is simply $\dbinom{(n-2)\mu + 3}{3}$. Taking the limit as $\mu$ approaches infinity, we obtain our closed form to be $\lim_{\mu \to \infty} \frac{\dbinom{(n-2)\mu + 3}{3}}{ (n\mu + 1)^3/3!  }.$ Dividing leading coefficients, we obtain $\lim_{\mu \to \infty} \frac{\dbinom{(n-2)\mu + 3}{3}}{ (n\mu + 1)^3/3!  } =   \frac{\mu(n-2)^3}{\mu n^3} = \frac{(n-2)^3}{n^3}$. Solving the inequality $\frac{(n-2)^3}{n^3} > \frac{1}{2}$, we get the least value of $n$ as $n=\boxed{10}$.

 -fidgetboss\_4000

 


\end{enumerate}